\subsection{单环上的模}

引理 3. —— 设 A 是一个单环，S 是一个单 A-模。用 D 表示 S 的交换子的反体。每个 A-模都同构于形如 \(S \otimes_{D} V\) 的 A-模，其中 V 是体 D 上的一个左向量空间。

这由 VIII, 第 120 页的引理 1 与森田定理（VIII, 第 103 页）得出。

命题 2. —— 设 A 是一个单环，S 是一个单 A-模。

\begin{enumerate}
\def\labelenumi{\alph{enumi})}
\item
  每个 A-模都是投射的且是 S 型同型的，因而是半单的。若它的长度为 a，则它同构于 \(\mathrm{S}^{(a)}\)。
\item
  每个非零 A-模都是生成模。
\item
  两个 A-模同构当且仅当它们有相同的长度。
\end{enumerate}

用 D 表示 S 的交换子的反体。由 VIII, 第 120 页的引理 1，S 是一个可逆 (A,D)-双模。设 M 是一个 A-模；由引理 3，它同构于形如 \(S \otimes_{D} V\) 的模，其中 V 是体 D 上的一个左向量空间。

D-模 V 是投射的且是 \(D_{s}\) 型同型的；它是生成模当且仅当它不约化为 0。最后，\(S \otimes_{D} V\) 的长度等于 D-向量空间 V 的维数，而两个向量空间同构当且仅当它们有相同的维数。于是命题 2 鉴于 VIII, 第 109 页的命题 10 与 VIII, 第 110 页的命题 12 得出。

设 \(r \geqslant 1\) 是一个整数。我们称基数 \(\mathfrak{a}\) 可被 \(r\) 整除，若存在一个基数 \(\mathfrak{b}\) 使 \(\mathfrak{a} = r\mathfrak{b}\)。若 \(\mathfrak{a}\) 是无限的，则情形如此，因为我们有 \(r\mathfrak{a} = \mathfrak{a}\)（《集合论》III, §6, No.~3, 第 188 页，推论 3）。由这一注记推出，若基数 \(\mathfrak{a}\) 可被 \(r\) 整除，则存在唯一的基数 \(\mathfrak{b}\) 使 \(\mathfrak{a} = r\mathfrak{b}\)。

推论. —— 设 k 是一个交换域（域），A 是 k 上一个有限次数的单 k-代数。每个单 A-模在 k 上都是有限维的。两个 A-模同构当且仅当它们在 k 上的维数相等。

每个单 A-模都同构于 \(A_{s}\) 的一个商，因此在 k 上是有限维的。于是该推论由命题 2, c) 得出。

命题 3. —— 设 A 是一个单环。一个 A-模 M 是自由的当且仅当它的长度可被 A 的长度整除。若情形如此，则 M 的所有基都有相同的基数，记作 \(\dim_{\mathrm{A}}(\mathrm{M})\)（II, §7, No.~3, 第 294 页，注 2），它由关系式

\[
\mathrm{long} _ {\mathrm{A}} (\mathrm{M}) = \mathrm{long} (\mathrm{A}) \cdot \mathrm{dim} _ {\mathrm{A}} (\mathrm{M}).\tag{1}
\]

设 M 是自由的，并设 \((e_{i})_{i\in\mathrm{I}}\) 是 M 的一个基。A-模 M 是诸 A-模 \(Ae_{i}\) 的直和，它们各自同构于 \(A_{s}\)。令 \(r=\operatorname{long}_{\mathrm{A}}(\mathrm{A}_{s})\)；这是一个大于或等于 1 的整数（VIII, 第 119 页，命题 1）。由 VIII, 第 73 页的公式 (13)，我们有 \(\operatorname{long}_{\mathrm{A}}(\mathrm{M})=r\operatorname{Card}(\mathrm{I})\)。

反之，设基数 \(\operatorname{long}_{\mathrm{A}}(\mathrm{M})\) 可被 r 整除。设 a 是使 \(\operatorname{long}_{\mathrm{A}}(\mathrm{M}) = r\mathfrak{a}\) 的基数。那么 A-模 M 与 \(\mathrm{A}_{s}^{(\mathfrak{a})}\) 有相同的长度，因而由命题 2 与之同构。这就证明 M 是自由的。

命题 4. —— 设 A 是一个单环，M 是一个非零 A-模。用 B 表示 A-模 M 的自同态环，并把 M 视为一个左 B-模。

\begin{enumerate}
\def\labelenumi{\alph{enumi})}
\item
  映射 \(a \mapsto a_{M}\) 是从 A 到 B-模 M 的自同态环的同构。
\item
  设 M 作为 A-模具有有限长度。那么环 B 是单环，并且我们有
\end{enumerate}

\[
\mathrm{long} _ {\mathrm{A}} (\mathrm{M}) = \mathrm{long} (\mathrm{B}) \text{ 且 } \mathrm{long} _ {\mathrm{B}} (\mathrm{M}) = \mathrm{long} (\mathrm{A}).\tag{2}
\]

由 VIII, 第 122 页的命题 2，A-模 M 是生成模。按定义我们有 \(\mathrm{B} = \mathrm{A}_{\mathrm{M}}'\)，于是断言 a) 由 VIII, 第 82 页的定理 2 得出。

设 M 是一个有限长度的 A-模。选取一个单 A-模 S，并设 D 是其自同态环的反体。由引理 3，A-模 M 同构于形如 \(S \otimes_{D} V\) 的模，其中 V 是体 D 上的一个左向量空间。向量空间 V 是有限维的。由 VIII, 第 64 页的定理 3，环 B 同构于 \(\operatorname{End}_{\mathrm{D}}(\mathrm{V})\)；由引理 2（VIII, 第 121 页），环 B 因而是单环。考虑到 VIII, 第 63 页的注 1，我们得到等式

\[
\mathrm{long(B)} = \mathrm {dim_ {D} (V)} = \mathrm {long_ {A} (M)}.
\]

由 VIII, 第 63 页的注 1 与 VIII, 第 121 页的注 3，我们有关系式

\[
\mathrm{long} _ {\mathrm{B}} (\mathrm{M}) = \mathrm{long} _ {\mathrm{End} _ {\mathrm{D}} (\mathrm{V})} (\mathrm{S} \otimes_ {\mathrm{D}} \mathrm{V}) = \dim_ {\mathrm{D}} (\mathrm{S}) = \mathrm{long} (\mathrm{A}),
\]

由此得最后一个关系式。

